\section{Methodology}
\label{sec:method}

We propose a formal framework for understanding LLMs as simulators based on manifold projection. We then describe the \method framework used to empirically validate this theory.

\subsection{Formal Framework: Simulation as Manifold Projection}
Let $\gM$ be a high-dimensional manifold representing the probability distribution of all human-documented behavior present in the LLM's training corpus. A specific simulation is defined by a conditioning context $c = (S, P)$, where $S$ is a social scenario and $P$ is a persona. The act of simulation can be seen as a projection $P_c: \gM \to \gP(O)$, where $\gP(O)$ is a probability distribution over the set of all possible behavioral outcomes $O$.

{\bf The Superset Hypothesis.} We hypothesize that for any scenario $S$, the distribution generated by the LLM, $P_{(S, P)}(O)$, can capture a "superset" of human outcomes by aggregating the behaviors of all personas $P \in \gP$. While a single human individual is limited to a single "branch" of behavior, the LLM maintains a superposition of these branches until sampled.

\subsection{\method Framework}
\begin{figure}[t]
    \centering
    \setlength{\unitlength}{1cm}
    \begin{picture}(12,4)
        \put(0,1.5){\framebox(2,1){Scenario $S$}}
        \put(0,0.5){\framebox(2,1){Persona $P$}}
        \put(2,1.5){\vector(1,0){1}}
        \put(3,1){\framebox(3,2){\gpt}}
        \put(6,2){\vector(2,1){2}}
        \put(6,2){\vector(2,-1){2}}
        \put(6,2){\vector(1,0){2}}
        \put(8.5,3){Outcome 1}
        \put(8.5,2){Outcome 2}
        \put(8.5,1){Outcome $N$}
        \put(3,0){\dashbox{0.1}(8,4){Multiverse Branching}}
    \end{picture}
    \caption{Overview of the \method framework. Given a scenario $S$ and a persona $P$, we roll out $N$ parallel continuations using \gpt to map the distribution of potential behavioral outcomes.}
    \label{fig:method}
\end{figure}

To explore this manifold, we employ a branching framework across three experimental dimensions.

{\bf Scenario Selection.} We use a "Mistaken Credit" social dilemma, where a character (the persona) is in a situation where they have been publicly given credit for a major achievement that was actually the work of a colleague. This scenario presents a high-stakes choice between social harmony, self-interest, and honesty.

{\bf Persona Conditioning.} We define four distinct personas to test the sensitivity of the outcome distribution to latent psychological factors:
\begin{itemize}[leftmargin=*,itemsep=0pt,topsep=0pt]
    \item \aggressive: High ambition, direct, willing to prioritize self-interest.
    \item \passive: High agreeableness, indirect, prioritizes social harmony.
    \item \strategic: Calculated, balances self-interest with strategic relationship management.
    \item \generic: A generic "human" prompt designed to elicit the model's base distribution.
\end{itemize}

{\bf Multiverse Sampling.} For each persona, we run $N=50$ parallel continuations using \gpt with a temperature of $T=0.7$. This allows us to observe the variance and entropy of the model's behavioral manifold for each persona.

\subsection{Evaluation Metrics}
To quantify the outcome manifold, we use a secondary \gptmini model to code each generated outcome based on the following behavioral factors:
\begin{itemize}[leftmargin=*,itemsep=0pt,topsep=0pt]
    \item {\bf Directness (1-10):} How explicitly the persona addresses the mistaken credit.
    \item {\bf Social Harmony (1-10):} How much the persona prioritizes the feelings of others and "face-saving."
    \item {\bf Behavioral Category:} A qualitative classification (e.g., "Collaborative Correction," "Immediate Confrontation," "Private Talk Later").
    \item {\bf Category Entropy:} We calculate the Shannon entropy $H = -\sum p_i \log p_i$ over the behavioral categories to measure the diversity of the outcome distribution.
\end{itemize}
